\documentclass[11pt,letterpaper]{article}

\usepackage[margin=1in]{geometry}
\usepackage[english]{babel}
\usepackage[utf8]{inputenc}
\usepackage{amsmath}
\usepackage{amssymb}
\usepackage{amsthm}
\usepackage{booktabs}
\usepackage{bm}
\usepackage{microtype}
\usepackage{mathtools}
\usepackage{natbib}
\usepackage[colorlinks=true,linkcolor=blue,citecolor=blue,urlcolor=blue]{hyperref}
\usepackage{tikz}
\usepackage{caption}
\usepackage{subfig}
\usepackage{graphicx}
\usepackage{float}
\usepackage{setspace}

\usetikzlibrary{arrows.meta, positioning, calc}

\theoremstyle{plain}
\newtheorem{theorem}{Theorem}
\newtheorem{proposition}{Proposition}
\newtheorem{lemma}{Lemma}
\newtheorem{corollary}{Corollary}
\newtheorem{conjecture}{Conjecture}

\theoremstyle{definition}
\newtheorem{definition}{Definition}
\newtheorem{assumption}{Assumption}
\newtheorem{example}{Example}

\theoremstyle{remark}
\newtheorem*{remark}{Remark}

\title{\bfseries Targeted Lobbying on Council Networks\\[0.4em]
{\large \emph{First Draft}}}
\author{Samir Kadariya \and Bryce Roberts\\[0.4em]
\small 14.18 Mathematical Economic Modeling \quad\textbar\quad MIT, Spring 2026}
\date{April 29, 2026}

\begin{document}

\maketitle

\begin{abstract}
\noindent A lobbyist with a limited budget must decide not only how hard to push a bill but \emph{whom} to push, since influence cascades through a council's network of party, committee, and personal ties. We model this as a static network game in the linear--quadratic tradition of \citet{ggg2020} and embed it in a voting institution that aggregates members' actions into a passage decision. Our main results are: (i) for a linear passage rule, the lobbyist optimally concentrates the entire budget on the member with the highest \emph{Katz--Bonacich centrality} with decay $\beta/c$, refining the eigenvector-centrality result of \citet{ggg2020}; and (ii) under any smooth strictly increasing vote-aggregation rule, the small-budget optimal target is the member with the largest \emph{slope-weighted Bonacich return} $w_j = \sum_i f'(a_i^0) M_{ij}$, where $a^0 = Mb$ is the no-lobbying equilibrium. When $f$ is linear or baseline actions are zero, $w$ collapses to Katz--Bonacich centrality; when baseline actions differ and $f$ is concave, the slope-weighting can overturn the Katz ranking, so a peripheral fence-sitter can be cheaper to target than a saturated central partisan. We illustrate both results on a three-member path network and discuss a probabilistic-vote companion model in which the swing-voter slope plays the same role as $f'(a^0)$.
\end{abstract}

\input{sections/01_intro}
\input{sections/02_related_work}
\input{sections/03_model_setup}
\input{sections/04_council_eq}
\input{sections/05_lobbyist_lp}
\input{sections/06_theorem2}
\input{sections/07_comparative_statics}
\input{sections/08_theorem3}
\input{sections/09_model_B}
\input{sections/10_discussion}
\input{sections/11_conclusion}

\bibliographystyle{plainnat}
\bibliography{refs}

\end{document}
